\section{Backward persistence and causal factorization}\label{sec:causality}

\begin{lemma}[Relative-boundary commutation]\label{lem:forward}
Two simultaneously legal formal bipyramid replacements with disjoint
consumed incarnation sets commute. Their two orders produce canonically
isomorphic marked triangulations, fixing all unaffected cells and respecting
the two replacement boundaries.
\end{lemma}
\begin{proof}
Replace both abstract fillings before imposing their attaching maps. Their
interiors are disjoint, and their individually labelled boundary triangles
are unchanged. Gluing the two new fillings to the old exterior therefore
gives the same quotient in either order. This includes a boundary triangle
of one region paired to a boundary triangle of the other.

The upward event's internal paired face and the downward event's complete
edge star of three tetrahedra are unaffected by a replacement with disjoint
consumed cells. Hence each event remains legal after the other. The five
formal height values are determined by the unchanged old edge differences;
the resulting rows agree after the tracked relabelling and normalization.
\end{proof}

The forward statement is insufficient to delete unrelated historical events:
a move might have become enabled only after an earlier move. For the present
model, such a dependency must involve an actual birth consumed by the later
event.

\begin{lemma}[Backward persistence]\label{lem:backward}
Suppose $a;b$ is a legal adjacent pair of moves and $b$ consumes no
incarnation born at $a$. Then $b$ was legal immediately before $a$, and
$a;b$ can be interchanged to $b;a$ under the tracked boundary identifications.
\end{lemma}
\begin{proof}
Every source incarnation of $b$ existed before $a$ and survived it. These
cells are disjoint from those consumed by $a$. Pairings between two survivor
tetrahedra, including their local permutations, are unchanged by $a$.

For an upward $b$, its internal face is exactly such a survivor-to-survivor
pairing. Its two distinct tetrahedra therefore support the same upward move
before $a$.

For a downward $b$, its three central local edge occurrences form a closed
triangular edge-link cycle. Every face incident to these local edges is
paired within the three survivor tetrahedra. Those same pairings existed
before $a$. Edge equivalence is generated by the pairings of faces containing
the local edge, so its orbit already closed within those three occurrences;
it could not have included a fourth exterior occurrence. The formal belt
and apex consistency checks involve only these unchanged permutations.
Thus the entire downward legality condition already held. Apply
\cref{lem:forward} to interchange the now simultaneously legal moves.
\end{proof}

\begin{definition}[Birth graph]
The birth graph $D_\sigma$ of a finite trace has its move events as vertices.
It has one directed arc $a\to b$ for each tetrahedron incarnation born at
$a$ and consumed at $b$. Parallel arcs are retained. The graph is acyclic,
and a component means a weakly connected component.
\end{definition}

\begin{theorem}[Component projection]\label{thm:projection}
The subsequence formed by any union of birth-graph components is a legal
trace from the original source, with the original local frames transported
through it. The full trace can be reordered into any order of its component
blocks. If a component $C$ has $u_C$ upward and $d_C$ downward events, then
the full loss is $\sum_C(d_C-u_C)$.
\end{theorem}
\begin{proof}
Preserve the internal order of the selected and unselected subsequences,
and move all selected events to the beginning by adjacent swaps. At a swap
$a;b$ with $a$ unselected and $b$ selected, there is no birth arc from $a$
to $b$: otherwise their components would coincide. Thus $b$ consumes no
output of $a$, and \cref{lem:backward} permits the interchange. Repeating
finitely many swaps makes the selected subsequence an initial legal trace.
The same procedure orders all component blocks. Marked endpoint agreement
comes from \cref{lem:forward}; the loss formula is \eqref{eq:balance}.
\end{proof}

This theorem preserves reachability of the projected trace, not every global
property of the full endpoint. Additivity of tetrahedron loss is essential
in the application below.

\section{Connected original supports and the incidence bound}

For a component $C$, let $F_C$ be the set of original tetrahedra consumed by
its events. The sets $F_C$ are disjoint, since an original incarnation is
consumed at most once. The following lemma concerns connectivity in the
\emph{original} dual graph, not just in some later triangulation.

\begin{lemma}[Original footprint connectivity]\label{lem:connected}
For every nonempty birth-graph component $C$, the induced graph
$G_T[F_C]$ is connected.
\end{lemma}
\begin{proof}
Maintain a partition of the original tetrahedra. Initially every block is
a singleton and labels its corresponding current incarnation. When a move
consumes incarnations labelled by blocks $A_1,\ldots,A_s$, merge these
blocks into $A=A_1\cup\cdots\cup A_s$. Relabel every surviving incarnation
of an affected block by $A$, and label all new incarnations by $A$.

Two invariants are maintained. First, each block induces a connected
subgraph of $G_T$. Second, between distinct blocks, current cut-face pairings
are in a boundary-port-preserving bijection with the original face pairings
crossing those blocks. Both assertions hold at the start.

The consumed tetrahedra of either move are connected in the current dual
graph: an upward move consumes an adjacent pair, and a downward move
consumes the triangular cycle around its central edge. Their distinct block
labels therefore form a connected graph. By the cut-face invariant, every
edge joining two such labels witnesses an original dual edge between the
corresponding original blocks. Their union is connected. The replacement
retains all six formal boundary triangles and their attaching permutations;
every pairing to an unaffected block keeps its original boundary port.
Pairings newly internal to the merged block cease to cross the partition.
This proves preservation of both invariants.

An activated block records exactly the original ancestors of a current weak
birth component. Consuming a born incarnation attaches its producing event
component; consuming births from several blocks merges those components;
an untouched original cell enters a block exactly at its first consumption.
Consequently the final activated blocks are precisely the $F_C$. The first
invariant proves the lemma.
\end{proof}

\begin{lemma}[Incidence bound]\label{lem:incidence}
A birth component with $u$ upward and $d$ downward events consumes at most
$u+2d+1$ original tetrahedra. More generally, a trace with $c$ birth
components satisfies
\begin{equation}\label{eq:incidence}
 |F(\sigma)|\le u+2d+c.
\end{equation}
\end{lemma}
\begin{proof}
Let $q$ count consumptions of born incarnations. Each consumed cell is
either original or born, and the move input sizes give the exact identity
\[
 |F(\sigma)|+q=2u+3d.
\]
The birth multigraph has $u+d$ vertices and $c$ weak components, hence at
least $u+d-c$ arcs. Its arc count is exactly $q$, including parallel arcs.
Subtract this lower bound from the identity. The connected case has $c=1$.
\end{proof}

\section{Automatic localization of a strict descent}\label{sec:localization}

\begin{theorem}[Causal localization]\label{thm:localization}
If $T$ admits any descending legal trace with at most $U$ total upward
moves, it admits a trace $\tau$ such that, for some $u\le U$:
\begin{align}
 u(\tau)&=u,& d(\tau)&=u+1,& |\tau|&=2u+1,\label{eq:counts}\\
 |T_\tau|&=t-1,& |F(\tau)|&\le3u+3\le3U+3.\label{eq:localized}
\end{align}
The birth graph of $\tau$ is connected, and so is its actual original
footprint in $G_T$.
\end{theorem}
\begin{proof}
Choose a minimum-length descending trace among those using at most $U$
upward moves. No earlier prefix descends. Every event changes the count
by exactly one, so the final count is exactly $t-1$, giving $d=u+1$.

If the birth graph had several components, their losses would sum to one.
Some component would have positive loss. Its projection is legal by
\cref{thm:projection}, uses no more upward moves, and is a strictly shorter
descending trace. This contradicts minimality. Hence the graph is connected.
By \cref{lem:connected}, its actual original footprint is connected; by
\cref{lem:incidence},
\[
 |F(\tau)|\le u+2(u+1)+1=3u+3.
\]
The length and remaining claims follow.
\end{proof}

\begin{remark}[Why one truncation is insufficient]
Extracting a descending component and truncating it at first descent does
not by itself prove that the retained birth graph remains connected. Later
events may have merged components that were previously separate. The
minimum-length argument obtains both properties at once. The implemented
certificate extractor instead alternates component projection and first-descent
truncation until stable; every nontrivial change removes an event.
\end{remark}

\begin{corollary}[Constructive certificate localization]\label{cor:extraction}
Given an explicit verified descending trace, a connected first-descending
subtrace satisfying \eqref{eq:counts}--\eqref{eq:localized} can be found and
replayed in time polynomial in the supplied trace's bit length.
\end{corollary}
\begin{proof}
Construct the birth graph by tracking live incarnations. Select a positive-loss
weak component and project to it, then truncate at its first strict descent.
Repeat. Each iteration retains a positive-loss trace and either terminates
with a connected first descent or removes at least one event. There are at
most as many iterations as input events, and each graph operation is polynomial.
Replay the retained events with their original local frames. The projection
theorem proves legality, and independent move and transport checkers verify
the emitted certificate. The incidence calculation then applies.
\end{proof}

\section{A sharp six-tetrahedron example}\label{sec:sharp}

Consider the simplicial ball with seven distinct vertices $A,B,C,D,E,X,Y$
and six tetrahedra
\begin{equation}\label{eq:gadget}
 ABCD,\ EBCD,\ ABCX,\ EBCX,\ ACDY,\ ECDY,
\end{equation}
gluing equally named triangular faces. Start with the bipyramid
$ABCD\cup EBCD$, then attach $ABCX\cup EBCX$ and $ACDY\cup ECDY$ along their
respective two-triangle boundary discs. Each attachment is a ball attached
along a disc, so the result is a $3$-ball.

\begin{figure}[ht]
\centering
\begin{tikzpicture}[every node/.style={font=\small},
  cell/.style={draw=inkblue,rounded corners=2pt,fill=lightblue,minimum width=19mm,minimum height=7mm}]
\node[cell] (a) at (0,0) {$ABCD$};
\node[cell] (b) at (0,-3.4) {$EBCD$};
\node[cell] (x) at (-3,-1.05) {$ABCX$};
\node[cell] (z) at (-3,-2.35) {$EBCX$};
\node[cell] (y) at (3,-1.05) {$ACDY$};
\node[cell] (w) at (3,-2.35) {$ECDY$};
\draw[inkblue,thick] (a)--node[right]{$BCD$}(b);
\draw[teal] (a)--node[above left]{$ABC$}(x);
\draw[teal] (x)--node[left]{$BCX$}(z);
\draw[teal] (z)--node[below left]{$EBC$}(b);
\draw[teal] (a)--node[above right]{$ACD$}(y);
\draw[teal] (y)--node[right]{$CDY$}(w);
\draw[teal] (w)--node[below right]{$ECD$}(b);
\end{tikzpicture}
\caption{The original dual graph of the sharp ball. Each labelled edge is
an internal paired triangle. The only internal primal edges are $BC$ and
$CD$, both of degree four. The central upward move enables two downward
moves, whose union consumes every original tetrahedron.}
\label{fig:gadget}
\end{figure}

\begin{proposition}[Sharpness at one upward move]\label{prop:sharp}
The ball \eqref{eq:gadget} has no initial downward move. It admits a descent
using one upward and two downward moves, but every descent using at most
one upward move consumes all six original tetrahedra. Thus the value
$3U+3=6$ cannot be reduced to five at $U=1$ in the stated move model.
\end{proposition}
\begin{proof}
The only internal edges are $BC$ and $CD$, each with four incident
tetrahedra. Therefore no downward move is initially legal. An upward move
through $BCD$ replaces the central two tetrahedra by
$AEBC,AECD,AEDB$. The edge $BC$ now supports a downward move using
$AEBC,ABCX,EBCX$, producing $AEBX,AECX$. The edge $CD$ supports the second
downward move using $AECD,ACDY,ECDY$, producing $AECY,AEDY$. Five tetrahedra
remain. The two downward moves commute in this displayed construction.

For the lower bound, a first descent with at most one upward move must have
the pattern up/down/down. The upward move creates one new internal edge.
If the first downward move deletes that edge, it undoes the upward move
and returns to a triangulation with no downward move. Otherwise it deletes
$BC$ or $CD$. To reduce that edge's degree from four to three, the upward
move must have used it in its central triangle. Its ensuing downward move
consumes exactly one of the three new tetrahedra and produces two tetrahedra
containing the newly created edge. That new edge's degree therefore changes
from three to four and it cannot be deleted by the second downward move.
The second move must delete the other of $BC,CD$.

An untouched original tetrahedron incident to an edge prevents that edge
from disappearing. Hence deleting both edges forces consumption of the
union of their original stars, which is precisely all six tetrahedra.
\end{proof}

The accompanying independent Regina inventory examines every first upward
move: seven candidates, fifteen up/down prefixes, and four successful
up/down/down traces. Their first faces are $BCD$ (two orders), $BCX$, and
$CDY$; all have footprint six. The initial and final isomorphism signatures
are respectively \code{gbbjrcadfaa} and \code{fbbfiaev}. These finite records
check the construction and the proof's edge inventory; they are not used to
infer an asymptotic statement from a few examples. The article claims sharpness
at $U=1$, not asymptotic optimality of the coefficient three for all $U$.

\section{A category boundary: strict simplicial moves}

The hypotheses about generalized triangulations cannot be omitted. Let $S$
be the boundary of an octahedron and form its simplicial suspension with
vertices $x,y$. This triangulates a $3$-sphere. For each base triangle $abc$,
the tetrahedra $xabc,yabc$ admit a strict simplicial $2$--$3$ move because
the proposed edge $xy$ is initially absent.

Choose two opposite base triangles. Their two move supports are disjoint,
and both moves are initially legal. Performing either creates the edge
$xy$, so the other is then forbidden by the strict simplicial nonexistence
condition for the new simplex. Disjoint consumed tetrahedra therefore do
not imply commutation in that category. Reversing the first move gives a
backward-persistence failure caused by deletion of a global obstruction,
without consumption of any tetrahedron used by the newly enabled move.

The repository's generalized model allows distinct edges with the same
global endpoints and has no such absent-simplex condition. Extending the
theorem to a different move model requires recording its additional
dependencies and reproving the locality bound.
